\paragraph{Modeling choice (smoothly bounded log-deviation).}
For each OD/time-bin cell, introduce an unconstrained latent variable
$\widetilde{z}_q^t\in\mathbb{R}$. To protect the exponentiation and the downstream
assignment from extreme values while retaining a smooth gradient, define the effective
log-deviation as
\[
z_q^t
=
z_{\max}\tanh\!\left(\frac{\widetilde{z}_q^t}{z_{\max}}\right),
\qquad z_{\max}>0.
\]
The demand used by the assignment is then
\[
f_q^t
=
f_{q}^{t,(0)}\exp(z_q^t).
\]
This transformation guarantees $f_q^t\ge 0$ whenever $f_{q}^{t,(0)}\ge 0$ and preserves the baseline
structure up to multiplicative adjustments. It also implies
\[
|z_q^t|\le z_{\max},
\qquad
\exp(-z_{\max})
\le \frac{f_q^t}{f_q^{t,(0)}}
\le \exp(z_{\max})
\]
for positive baseline cells. The default $z_{\max}=6$ therefore permits demand
multipliers ranging approximately from $1/403$ to $403$. In the software configuration,
this smooth bound is exposed under the historical name \texttt{z\_clip}; despite that
name, the implementation uses a smooth hyperbolic-tangent transformation rather than
hard clipping.

